% main.tex - 整书主控模板
%
% 作用：整书合并时所有章节的唯一入口。把各章 \input 进来，统一页面、页眉页脚、
% 参考文献样式、代码块、图片宽度上限等。章节源文件无需关心这些全局设置。
%
% 使用步骤：
%   1. 替换所有 {{XXX}} 占位符（TITLE / SUBTITLE / EPIGRAPH / EPIGRAPH_CLOSING / AUTHOR / DATE / PDF_TITLE）。
%      - {{TITLE}} 用于扉页显示，可含 \\ 换行；{{PDF_TITLE}} 用于 PDF 元数据，必须单行、不含 \\。
%      - {{EPIGRAPH_CLOSING}} 用于结尾段引文（可选，留空则整段删除）。
%      - 注：共用补丁占位符（SHARED_PREAMBLE）不是给用户填的——scaffold 会自动替换为
%        共用补丁片段（中文 URL 支持 / 参考文献降级 / codeblock 环境 / 防溢出），无需手动改动，
%        其唯一真源是 templates/_preamble_shared.tex。
%   2. 在下方“正文各章”处的章节列表占位符中，每章写一行 \input{第N章/第N章}。
%   3. 若不需要扉页 / 序言 / 附录 / 结尾段，删除对应整块。
%   4. xelatex 跑两遍（第二遍解析交叉引用）。
%   5. 收尾：删 main.aux/log/out/toc，保留 main.tex 与 main.pdf。

\documentclass[a4paper, 11pt, twoside]{ctexrep}
\usepackage{amsmath,amssymb}
\usepackage{graphicx}
\usepackage{booktabs}
\usepackage{multirow}
\usepackage{array}
\usepackage{longtable}
\usepackage{tabularx}
\usepackage[unicode]{hyperref}
\usepackage{indentfirst}
\usepackage{xcolor}
\usepackage{xeCJK}
\usepackage{tikz}
\usetikzlibrary{shapes.geometric, arrows.meta, positioning, fit, backgrounds, calc}
\usepackage{adjustbox}            % 图片宽度上限：每个 tikzpicture 套 max width=\textwidth
\usepackage{tcolorbox}
\tcbuselibrary{listings, skins, breakable}   % listings 引擎让 codeblock 自动折行
\usepackage{enumitem}
\usepackage{url}

\usepackage{etoolbox}
\usepackage{geometry}

% === 共用补丁片段（唯一真源：templates/_preamble_shared.tex） ===
% scaffold 会把下方独立成行的占位符替换为本片段内容（中文 URL 支持 / 参考文献降级 /
% codeblock 环境 / 防溢出），避免与 _tmp.tex 双份手维护导致漂移。
% 此时 etoolbox（\patchcmd / \AtBeginEnvironment）、tcolorbox（\newtcblisting）
% 均已加载，5 段补丁的顺序依赖全部满足。
{{SHARED_PREAMBLE}}

\geometry{margin=2.5cm}

% === 页眉页脚统一 ===
\setlength{\headheight}{15pt}
\makeatletter
\newcommand{\pdfm@noUppercase}[2][]{#2}
\long\def\pdfm@hdrshow#1{{%
  \let\uppercase\relax
  \let\MakeUppercase\pdfm@noUppercase
  \expandafter\let\csname MakeUppercase \endcsname\relax
  \expandafter\def\csname MakeUppercase\space\space\space\endcsname[##1]##2{##2}%
  #1}}
\def\ps@pdfmaker{%
  \def\@evenhead{\pdfm@hdrshow{\slshape\small\leftmark}\hfil}%
  \def\@oddhead{\hfil\pdfm@hdrshow{\slshape\small\rightmark}}%
  \def\@evenfoot{\hfil\thepage\hfil}%
  \def\@oddfoot{\hfil\thepage\hfil}%
}
\makeatother
\pagestyle{pdfmaker}

\XeTeXlinebreaklocale "zh"
\XeTeXlinebreakskip=0pt
\graphicspath{{figures/}}
\DeclareGraphicsExtensions{.pdf,.png,.jpg}

\hypersetup{
  colorlinks=true, linkcolor=blue, urlcolor=blue, citecolor=blue,
  pdftitle={{PDF_TITLE}},
  pdfauthor={{AUTHOR}},
  pdfsubject={{SUBTITLE}}
}

\begin{document}

% ============ 扉页 ============
\thispagestyle{empty}
\begin{center}
\vspace*{5cm}
{\Huge\bfseries {{TITLE}}\par}
\vspace{0.8cm}
{\Large\bfseries {{SUBTITLE}}\par}
\vspace{1.5cm}
\rule{0.7\textwidth}{0.4pt}\par
\vspace{1cm}
{\large {{EPIGRAPH}}\par}
\vspace{4cm}
{\large {{AUTHOR}}\quad 整理\par}
\vspace{0.8cm}
{\large {{DATE}}\par}
\end{center}
\newpage

% ============ 摘要（必选：build 合并前会预检其存在性） ============
\thispagestyle{empty}
\chapter*{摘\quad 要}
\addcontentsline{toc}{chapter}{摘要}

% 在此输入摘要正文（建议 150-300 字，概括全书主题、章节结构与核心结论）

\newpage

% ============ 序言（可选：若不需要则整段删除） ============
\thispagestyle{empty}
\chapter*{序\quad 言}
\addcontentsline{toc}{chapter}{序言}

% 在此输入序言正文（建议 200-400 字，介绍写作缘起 / 资料来源 / 读者对象 / 阅读建议）

% ⚠️ 若在前置部分（序言/术语表/摘要）插入整宽表格（tabularx/tabular/table），
% 务必在其前一行加 \noindent：否则段落缩进（\parindent，11pt 下 2em≈21.9pt）
% 会与整宽表格叠加，触发恰好 21.9pt 的 Overfull（build 合并前会专门预检此项）。
% 示例：
%   \noindent\begin{tabularx}{\textwidth}{l X} ... \end{tabularx}

\newpage

% ============ 目录 ============
\thispagestyle{empty}
\renewcommand{\contentsname}{目\quad 录}
\tableofcontents
\newpage

% ============ 正文各章 ============
% 章节文件用 \chapter{} 作为顶层命令（与 ctexrep 默认 chapter 行为对齐）。
% 每章一个 \input；ctexrep 自动生成「第 X 章」前缀与 X.Y 编号。
{{CHAPTERS_LIST}}

% ============ 附录（可选：若不需要则整段删除） ============
% \appendix
% {{APPENDICES_LIST}}

% ============ 结尾段（可选：若不需要则整段删除） ============
\vspace{2em}
\begin{center}
\rule{0.5\textwidth}{0.4pt}\par
\vspace{0.5cm}
{\large\itshape {{EPIGRAPH_CLOSING}}\par}
\vspace{0.3cm}
{\small {{DATE}}\ 完稿\par}
\end{center}

\end{document}
